% /Users/pikachu/books/payments-book/src/parts/part5/ch28-processing-center.tex
\chapter{Процессинговый центр}
\label{ch:processing-center}

\begin{kaobox}[title=Что вы узнаете из этой главы]
  Что стоит за~\enquote{чёрным ящиком}, который за~сотню
  миллисекунд отвечает на~авторизационный запрос: из~каких
  компонентов состоит процессинговый центр, как~распределяется
  latency budget между ними, почему active-active~--- не~роскошь,
  а~необходимость, как~центр подключается к~карточным сетям
  и~какие лицензии для~этого нужны.
\end{kaobox}

Процессинговый центр удобно мыслить как машину, одновременно
подчинённую трём инвариантам: ответить быстро, не~остановиться
при~отказе узла и~не~выпустить ключевой материал за~пределы
защищённого контура. Все его компоненты и~все архитектурные
компромиссы дальше будут крутиться вокруг latency, availability
и~custody of keys.

\section{Компоненты процессингового центра}
\label{sec:pc-components}

Платёжный шлюз, описанный в~предыдущей главе, принимает запрос
мерчанта и~транслирует его в~сообщение \iso{8583}. Но~между
шлюзом и~карточной сетью стоит ещё одна система, без~которой
ни~авторизация, ни~клиринг, ни~settlement невозможны,~---
процессинговый центр.

\begin{termdef}{Процессинговый центр}
  Процессинговый центр (processing center)~--- комплекс
  программно-аппаратных средств, выполняющий авторизацию
  карточных транзакций, формирование клиринговых файлов,
  расчёт settlement-позиций и~криптографические операции
  (верификация PIN, генерация и~проверка криптограмм).
  В~российском праве процессинговый центр может выступать
  как~операционный и~платёжный клиринговый центр (ОПКЦ),
  если имеет соответствующую лицензию ЦБ.
\end{termdef}

Внутри процессингового центра можно выделить четыре основных
компонента, каждый из~которых решает свою задачу и~предъявляет
свои требования к~производительности и~отказоустойчивости.

\subsection*{Authorization switch}

Authorization switch~--- ядро процессинга, принимающее
авторизационные сообщения \iso{8583} от~эквайеров (или~от
шлюзов, работающих от~имени эквайеров), маршрутизирующее их
к~нужному эмитенту (или~в~карточную сеть, которая доставит
сообщение эмитенту) и~возвращающее ответ. Switch работает
в~режиме реального времени: каждое сообщение обрабатывается
индивидуально, а~допустимая латентность измеряется десятками
миллисекунд.

Маршрутизация внутри switch определяется BIN-таблицей:
для~каждого BIN-диапазона известно, через какой канал
(VisaNet, Banknet, ОПКЦ НСПК, прямое подключение к~эмитенту)
следует отправить запрос. BIN-таблица обновляется регулярно~---
карточные сети публикуют обновления BIN-диапазонов, и~процессинг
обязан применять их в~установленные сроки, иначе транзакции
будут маршрутизироваться неверно и~отклоняться.

Помимо маршрутизации, switch выполняет базовую валидацию
сообщения (проверка обязательных полей, формат MTI, корректность
bitmap), обогащение (добавление полей, требуемых конкретной
сетью, но~отсутствующих в~исходном запросе) и~протоколирование:
каждая транзакция записывается в~transaction log с~точностью
до~миллисекунд, что~необходимо как~для~отладки, так и~для
последующей сверки с~клиринговыми файлами.

Отдельная функция switch~--- stand-in processing: если~карточная
сеть или~эмитент недоступны (timeout, сетевой сбой), switch
может принять решение об~авторизации самостоятельно, опираясь
на~локальные правила (лимит суммы, BIN-range, история транзакций
по~карте). Stand-in~--- вынужденная мера: процессинг берёт
на~себя кредитный риск, который в~нормальном режиме несёт
эмитент, и~поэтому stand-in ограничен низкими суммами
и~определёнными категориями транзакций.

\subsection*{HSM farm}

Криптографические операции в~процессинговом центре выполняются
не~программно, а~на~специализированных аппаратных модулях~---
HSM (Hardware Security Module), физически защищённых устройствах,
хранящих ключи в~tamper-resistant среде и~выполняющих
криптографические функции внутри защищённого
периметра~(см.~главу~\ref{ch:crypto}).

HSM farm в~процессинге решает несколько задач. При~авторизации~---
верификация PIN: switch извлекает PIN block из~\de{52}, отправляет
его в~HSM вместе с~ключом зоны (Zone PIN Key), и~HSM
транслирует PIN block из~формата эквайера в~формат эмитента,
используя соответствующие ключи. При~EMV-транзакции~---
проверка криптограммы ARQC: HSM получает данные транзакции
и~ключ эмитента (или~его~дериват) и~вычисляет ожидаемую
криптограмму, сравнивая её~с~полученной от~карты. При~генерации
ответа эмитента~--- формирование ARPC (Authorization Response
Cryptogram), которая передаётся карте для~взаимной
аутентификации.

HSM~--- узкое место по~латентности: каждая криптографическая
операция занимает 1--5~мс, и~при~нагрузке в~тысячи транзакций
в~секунду одного HSM не~хватает. Поэтому процессинговые центры
используют HSM farm~--- пул из~десятков устройств
с~балансировкой нагрузки, где~каждое устройство содержит
идентичный набор ключей, загруженных при~инициализации
(key injection). Отказ одного HSM не~должен прерывать
обработку: оставшиеся устройства принимают нагрузку,
а~замена происходит без~остановки сервиса.

\begin{important}
  HSM~--- единственный компонент процессинга, для~которого
  горизонтальное масштабирование ограничено физически: каждое
  новое устройство требует церемонии загрузки ключей
  (key ceremony), которая проводится в~защищённом помещении
  с~участием нескольких custodians и~занимает часы. Планирование
  ёмкости HSM farm должно учитывать не~только текущую нагрузку,
  но~и~время, необходимое для~ввода в~эксплуатацию нового
  устройства.
\end{important}

\subsection*{Clearing engine}

Если authorization switch работает в~реальном времени,
то~clearing engine~--- пакетный (batch) процесс, запускаемый
по~расписанию (обычно один-два раза в~сутки). Его~задача~---
собрать все авторизованные транзакции за~период, обогатить их
данными, которых не~было на~этапе авторизации (финальная сумма
после tip adjustment, результат адресной верификации,
DCC-конверсия), и~сформировать клиринговый файл в~формате,
требуемом карточной сетью: TC33 для~\scheme{Visa},
IPM для~\scheme{Mastercard}, собственный формат
для~НСПК~(см.~главу~\ref{ch:reconciliation}).

Clearing engine сопоставляет каждую клиринговую запись
с~авторизацией по~ключевым полям (STAN, RRN, сумма, дата)
и~выявляет расхождения: авторизация без~клиринга (orphan
authorization), клиринг без~авторизации (late presentment,
force post), несовпадение сумм (partial capture, tip
adjustment). Каждое расхождение~--- потенциальная финансовая
потеря, и~clearing engine генерирует exception report,
который операционная команда разбирает вручную или~по~правилам
автоматической обработки.

\subsection*{Settlement module}

Settlement module рассчитывает нетто-позиции участников:
сколько каждый эмитент должен заплатить и~сколько каждый
эквайер должен получить, с~учётом interchange, комиссий
сети и~корректировок (chargebacks, representments, fees).
Результат~--- набор проводок, которые инициируются через
расчётную инфраструктуру: для~внутрироссийских транзакций~---
через НСПК и~платёжную систему Банка
России~(см.~главу~\ref{ch:banks}), для~международных~---
через расчётные банки карточных сетей.

Settlement~--- финальная стадия, после которой деньги фактически
перемещаются между банками; до~этого момента все~операции~---
лишь записи в~учётных системах процессинга. Ошибка
в~settlement~--- это~реальная финансовая потеря, а~не~
техническая проблема, которую можно исправить retry.


Итак, карта компонентов уже показывает, что процессинг~--- это
не~один switch, а~связка real-time маршрутизации, пакетной
обработки и~аппаратной криптографии. Следующий вопрос естественен:
как уложить такую систему в~жёсткий latency budget и~одновременно
сделать её отказоустойчивой.

\section{Производительность и отказоустойчивость}
\label{sec:pc-performance}

Процессинговый центр~--- одна из~немногих систем, где~требования
к~латентности и~доступности не~маркетинговое украшение,
а~контрактное обязательство: карточные сети устанавливают SLA
на~время ответа и~штрафуют за~его~нарушение.

\subsection*{Latency budget}

Типичное требование карточной сети~--- ответ на~авторизационный
запрос в~течение 3--5~секунд от~момента получения сообщения
до~отправки ответа. Но~3--5~секунд~--- это~максимум, включающий
время на~доставку сообщения эмитенту и~обратно; на~долю
процессинга остаётся значительно меньше. Хорошо спроектированный
процессинговый центр укладывается в~p99~$<$~100~мс
на~обработку одного авторизационного сообщения, и~этот бюджет
распределяется примерно так:

\begin{itemize}
  \item Приём и~парсинг сообщения \iso{8583}: 1--2~мс.
  \item BIN-lookup и~маршрутизация: 1--3~мс.
  \item Обращение к~HSM (PIN translation или~проверка
        криптограммы): 3--10~мс.
  \item Запись в~transaction log: 1--2~мс.
  \item Сериализация ответа и~отправка: 1--2~мс.
  \item Сетевая латентность до~карточной сети: 10--50~мс
        (зависит от~географии).
\end{itemize}

Суммарно внутренняя обработка занимает 10--20~мс; основная
часть бюджета уходит на~ожидание ответа от~эмитента (через
карточную сеть), на~которое процессинг повлиять не~может.
Именно поэтому оптимизация внутренней обработки, хотя
и~необходима, даёт ограниченный эффект: разница между
10 и~20~мс незаметна на~фоне 500--2000~мс ожидания ответа
от~удалённого эмитента.

\subsection*{Active-active и geo-redundancy}

Процессинговый центр, обрабатывающий карточные транзакции,
не~может позволить себе плановый простой: даже несколько минут
недоступности~--- это~тысячи отклонённых транзакций, недовольные
держатели карт и~штрафы от~карточных сетей. Поэтому стандартная
архитектура~--- active-active с~двумя (или~более) ЦОД,
каждый из~которых способен обрабатывать полный объём трафика.

В~active-active конфигурации трафик распределяется между ЦОД
в~нормальном режиме (обычно~50/50 или~по~географическому
признаку), и~при~отказе одного весь трафик переключается
на~оставшийся. Ключевая сложность~--- синхронизация состояния
между ЦОД: transaction log, BIN-таблицы, stand-in лимиты,
HSM-ключи должны быть идентичны в~обоих центрах.

Для~авторизации, которая по~природе stateless (каждое сообщение
обрабатывается независимо), синхронизация минимальна:
достаточно, чтобы BIN-таблицы и~ключи были одинаковы, а~transaction
log реплицировался асинхронно. Для~клиринга и~settlement,
которые работают с~накопленным состоянием, требования строже:
потеря клиринговых данных означает потерю денег, и~RPO
(Recovery Point Objective) должен быть близок к~нулю.

\begin{important}
  Для~авторизации RPO может быть несколько секунд: потеря
  нескольких транзакций из~transaction log~--- неприятность,
  но~не~катастрофа, потому что~авторизации можно восстановить
  из~логов карточной сети. Для~клиринга и~settlement RPO
  должен быть нулевым: потерянная клиринговая запись~---
  это~транзакция, за~которую эквайер не~получит деньги
  или~эмитент не~спишет средства. Эта~асимметрия определяет
  архитектуру репликации: авторизация допускает асинхронную
  репликацию ради~минимальной латентности, а~клиринг требует
  синхронной~--- ради~консистентности.
\end{important}

\begin{practice}
  Шардинг по~BIN-диапазонам~--- простейшая и~наиболее
  распространённая стратегия горизонтального масштабирования
  authorization switch. BIN однозначно определяет эмитента
  и~карточную сеть, а~значит~--- маршрут и~набор ключей;
  транзакции с~разными BIN не~зависят друг от~друга
  и~могут обрабатываться разными инстансами switch
  без~координации. При~необходимости шард можно перенести
  на~другой сервер или~в~другой ЦОД без~влияния на~остальные
  BIN-диапазоны.
\end{practice}

\subsection*{Disaster Recovery}

Geo-redundancy~--- расположение ЦОД в~разных географических
точках~--- защищает от~региональных катастроф (пожар, наводнение,
отключение электричества в~районе). Расстояние между ЦОД~---
компромисс: чем~дальше друг от~друга, тем~лучше защита
от~региональных событий, но~тем~выше сетевая латентность
синхронной репликации. Типичное расстояние~--- 50--200~км:
достаточно для~защиты от~локальных инцидентов при~сетевой
задержке менее~5~мс.

RTO (Recovery Time Objective)~--- время восстановления после
полного отказа ЦОД~--- для~процессинга измеряется секундами,
а~не~минутами: переключение трафика на~резервный ЦОД должно
происходить автоматически, без~ручного вмешательства,
по~сигналу от~системы мониторинга. На~практике RTO~$<$~30~секунд~---
ожидаемый уровень для~tier-1 процессингового центра.


\section{Подключение к~платёжным системам}
\label{sec:pc-connectivity}

Процессинговый центр~--- не~изолированная система: он~существует
в~сети подключений к~карточным системам, каждое из~которых
требует сертификации, выделенных каналов связи и~соблюдения
специфических протоколов.

\subsection*{VisaNet}

Подключение к~VisaNet~--- глобальной сети \scheme{Visa}~---
осуществляется через выделенные каналы (leased lines
или~VPN поверх интернета с~шифрованием). Авторизация проходит
через BASE~I~--- систему реального времени, обрабатывающую
авторизационные запросы; клиринг~--- через BASE~II, принимающую
batch-файлы. Процессинг, подключающийся к~VisaNet, проходит
многоэтапную сертификацию: тестирование форматов сообщений,
криптографии, stand-in логики, обработки reversal
и~timeout~\todocite{Visa Acceptance Testing Guide~---
  описание процедуры сертификации}.

\subsection*{Banknet}

\scheme{Mastercard} использует собственную сеть~--- Banknet,
архитектурно аналогичную VisaNet, но~с~рядом протокольных
отличий. Ключевое из~них~--- DMSA (Dual Message System
Architecture), при~которой \scheme{Mastercard} допускает как~dual
message (авторизация отдельно, клиринг отдельно), так и~single
message (авторизация и~клиринг в~одном сообщении)~---
в~отличие от~\scheme{Visa}, исторически ориентированной
на~dual message. Клиринговые файлы \scheme{Mastercard}
передаются в~формате IPM через систему
GCMS~\todocite{Mastercard Interface Processing Manual~---
  формат IPM, структура GCMS}.

\subsection*{ОПКЦ НСПК}

Все внутрироссийские карточные транзакции~--- независимо
от~того, выпущена карта \scheme{Visa}, \scheme{Mastercard}
или~Мир,~--- маршрутизируются через операционный и~платёжный
клиринговый центр (ОПКЦ) НСПК. Это~означает, что~процессинговый
центр российского банка подключается не~напрямую к~VisaNet
или~Banknet для~внутренних транзакций, а~к~ОПКЦ, который
выступает единым хабом: принимает авторизационные запросы
от~эквайеров, маршрутизирует их~к~эмитентам и~обеспечивает
клиринг и~settlement~(см.~главу~\ref{ch:mir}).

Подключение к~ОПКЦ требует сертификации по~стандартам НСПК
и~поддержки специфических расширений \iso{8583}, используемых
в~российской платёжной инфраструктуре~--- в~частности, полей
для~передачи данных платёжного приложения Мир (МРА)
с~российской криптографией (ГОСТ).

\subsection*{СБП}

Подключение к~СБП архитектурно отличается от~карточного
процессинга: если~карточные транзакции используют \iso{8583}
и~двухфазную модель (авторизация~+~клиринг), то~СБП работает
по~API-based протоколу с~финальными платежами~--- нет~авторизации
как~отдельной фазы, нет~chargeback, и~расчёт происходит
в~реальном времени~(см.~главу~\ref{ch:sbp}). Процессинговый
центр, поддерживающий и~карты, и~СБП, вынужден совмещать
два принципиально разных flow в~одной инфраструктуре.

\subsection*{Backpressure и graceful degradation}

Что~происходит, когда карточная сеть не~отвечает? Процессинговый
центр не~может просто ждать: тысячи новых авторизаций продолжают
поступать, и~если каждая из~них будет ждать ответа от~неотвечающей
сети, очередь переполнится и~парализует весь процессинг.

Грамотная реализация включает несколько уровней защиты.
Первый~--- timeout на~уровне соединения с~сетью (обычно
30--60~секунд), после которого транзакция либо~отклоняется,
либо~переводится в~stand-in. Второй~--- circuit breaker:
если доля timeout превышает порог (например, 50\% за~последнюю
минуту), процессинг прекращает отправку в~эту~сеть и~переключает
все транзакции на~stand-in или~альтернативный маршрут.
Третий~--- backpressure к~эквайерам: если процессинг перегружен,
он~начинает отвечать \rc{96}~(System Malfunction), сигнализируя
эквайерам о~временной недоступности, чтобы те~могли задействовать
собственную логику failover.


\section{Лицензирование и сертификация}
\label{sec:pc-licensing}

Процессинговый центр работает на~пересечении требований
регулятора, карточных сетей и~индустриальных стандартов
безопасности, и~каждый из~этих уровней предъявляет свои
требования.

\subsection*{Требования Банка России}

В~российском праве процессинговый центр может выступать
как~оператор услуг платёжной инфраструктуры (ОУПИ),
а~также как~операционный центр и~платёжный клиринговый центр.
Для~каждой роли предусмотрены свои лицензионные требования
в~рамках \law{161-ФЗ}: минимальный капитал, требования
к~бесперебойности, обязательства по~предоставлению отчётности
в~ЦБ~(см.~главу~\ref{ch:regulation}).

Кроме того, процессинг обязан выполнять требования
\law{382-П} (защита информации при~осуществлении переводов
денежных средств) и~\law{683-П} (требования к~защите
информации в~кредитных организациях), которые устанавливают
конкретные технические меры~--- от~шифрования каналов связи
до~контроля доступа и~логирования операций.

\subsection*{Сертификация у~карточных сетей}

Каждая карточная сеть проводит собственную программу
сертификации для~процессинговых центров, подключающихся
к~её~инфраструктуре. Сертификация включает:

\begin{itemize}
  \item Тестирование форматов сообщений: корректность заполнения
        data elements \iso{8583} для~всех типов транзакций
        (авторизация, reversal, клиринг, chargeback).
  \item Тестирование криптографии: корректность PIN translation,
        проверки EMV-криптограмм, работы с~ключами.
  \item Тестирование отказоустойчивости: обработка timeout,
        stand-in, duplicate detection.
  \item Compliance review: соответствие правилам сети (Core Rules
        для~\scheme{Visa}, Network Rules
        для~\scheme{Mastercard}).
\end{itemize}

Сертификация~--- не~разовое событие: при~существенных изменениях
в~протоколах (новые MTI, новые data elements, изменения
в~криптографии) карточные сети требуют ресертификации,
и~процессинг обязан пройти её~в~установленные сроки.

\subsection*{PCI DSS}

Процессинговый центр по~определению входит в~CDE (cardholder
data environment): через него~проходят PAN, PIN block,
CVV, криптограммы~--- полный набор данных держателя
карты~(см.~главу~\ref{ch:pci-dss}). Это~означает PCI~DSS
Level~1~--- максимальный уровень, требующий ежегодного
аудита квалифицированным QSA (Qualified Security Assessor)
и~квартальных ASV-сканов (Approved Scanning Vendor).

Для~процессинга, хранящего и~обрабатывающего PIN, дополнительно
применяется PCI PIN Security~--- отдельный стандарт, регулирующий
физическую и~логическую защиту PIN, HSM-процедуры, key management
и~уничтожение ключей.


\section*{Итоги главы}

\begin{itemize}
  \item Процессинговый центр состоит из~четырёх компонентов:
        authorization switch (реальное время, маршрутизация
        по~BIN), HSM farm (криптография~--- PIN, ARQC, ARPC),
        clearing engine (batch-формирование TC33/IPM)
        и~settlement module (расчёт нетто-позиций).
  \item Latency budget авторизации: p99~$<$~100~мс на~внутреннюю
        обработку; основная доля задержки~--- ожидание ответа
        от~эмитента через карточную сеть, на~которое процессинг
        повлиять не~может.
  \item Active-active с~двумя ЦОД~--- стандартная архитектура;
        авторизация допускает асинхронную репликацию,
        клиринг требует синхронной; RTO~$<$~30~секунд;
        шардинг по~BIN~--- простейшая стратегия масштабирования.
  \item Процессинг подключается к~VisaNet (BASE~I/II),
        Banknet (DMSA, IPM), ОПКЦ НСПК (единый хаб
        для~внутрироссийских транзакций) и~СБП (API-based,
        финальные платежи); при~недоступности сети~---
        stand-in, circuit breaker, backpressure (\rc{96}).
  \item Лицензирование включает три уровня: ЦБ~РФ
        (\law{161-ФЗ}, \law{382-П}), сертификация
        у~карточных сетей (форматы, криптография, compliance)
        и~PCI~DSS Level~1 + PCI PIN Security.
\end{itemize}
